\section{Discussion}
\label{sec:discussion}

% =============================================================================
% 5.1  Physical interpretation: why frequency drops with ripening
% =============================================================================
\subsection{Why the tap-tone drops with ripening}
\label{sec:physical_interp}

The monotonic decrease in tap-tone frequency during ripening
(Fig.~\ref{fig:ripening_sweep}) has a straightforward physical
explanation.  As the fruit ripens, enzymatic degradation of pectin in the
rind cell walls reduces the effective Young's modulus
$E_{\mathrm{rind}}$.  Since $f_2 \propto \sqrt{E_{\mathrm{rind}}}$
(Eq.~\ref{eq:f2_squared}), the frequency falls.  This is precisely the
information that experienced shoppers extract unconsciously: a lower pitch
signals a softer, riper melon.

The frequency shift is not driven by changes in fluid density or fruit
geometry, both of which vary by less than $\SI{5}{\percent}$ over the
ripening window.  The Sobol analysis (\S\ref{sec:sobol}) confirms that
$E_{\mathrm{rind}}$ is the dominant parameter, contributing more than
half the total variance.

% =============================================================================
% 5.2  Connection to biomedical FSI
% =============================================================================
\subsection{Connection to biomedical vibroacoustics}
\label{sec:biomedical}

The shell model employed here is structurally identical to the one
developed in our companion papers for the human abdomen.  Both systems
are fluid-filled, viscoelastic, oblate spheroidal shells; only the
material parameters differ.  The watermelon rind
($E \sim \SIrange{2}{20}{\mega\pascal}$) is considerably stiffer than
the abdominal wall ($E \sim \SI{0.1}{\mega\pascal}$), which shifts the
eigenfrequencies upward by roughly an order of magnitude---from
$\sim\SI{4}{\hertz}$ for the abdomen to $\sim\SIrange{80}{200}{\hertz}$
for watermelons.

This connection is more than a mathematical curiosity.  It demonstrates
that a single theoretical framework, rooted in Lamb's classical shell
theory~\cite{lamb1882} and the Junger--Feit fluid-loading
formulation~\cite{Junger1986}, can span applications from clinical
diagnostics to agricultural quality assessment.

% =============================================================================
% 5.3  Limitations
% =============================================================================
\subsection{Limitations}
\label{sec:limitations}

Several simplifications deserve acknowledgement:

\begin{itemize}
  \item \textbf{Isotropic rind.}  The watermelon rind is a biological
    composite with fibre orientation that may induce anisotropy.  We
    treat it as isotropic; orthotropic extensions are straightforward
    but require additional material characterisation.

  \item \textbf{Uniform thickness.}  Real rinds are thicker near the
    blossom end.  Thickness variation could be incorporated via
    perturbation methods but is unlikely to shift $f_2$ by more than
    $\sim\SI{10}{\percent}$.

  \item \textbf{Flesh as inviscid fluid.}  Watermelon flesh has finite
    shear strength.  At the frequencies of interest
    ($\sim\SI{100}{\hertz}$), the flesh is well within the
    fluid-like regime ($\omega \tau \ll 1$, where $\tau$ is the
    viscoelastic relaxation time), so this approximation is reasonable.

  \item \textbf{Turgor pressure.}  We neglect internal turgor pressure,
    which stiffens the shell.  For ripe fruit with degraded cell
    walls, this is appropriate; for very unripe fruit, turgor may
    contribute a correction of order $\SI{10}{\percent}$.
\end{itemize}

% =============================================================================
% 5.4  Practical implications
% =============================================================================
\subsection{Practical implications}
\label{sec:practical}

The closed-form inversion (Eq.~\ref{eq:inversion}) is computationally
trivial and could be implemented in a smartphone application.  A
microphone records the tap-tone, a fast Fourier transform extracts
$f_2$, and the inversion formula returns $E_{\mathrm{rind}}$ given
approximate fruit dimensions (obtainable from a photograph).  The
dimensionless ripeness parameter $\Pi_{\mathrm{ripe}}$ then maps the
result onto a cultivar-independent ripeness scale.

This approach offers several advantages over existing non-destructive
methods (near-infrared spectroscopy, impact testing with accelerometers):
it requires no specialised hardware, is non-contact (or nearly so---one
need only tap), and provides a physics-based rather than purely
statistical estimate of internal quality.
